\section{Режекторный полосовой фильтр}

Перейдем ко второму заданию, в рамках которого необходимо проанализировать работу режекторного фильтра. В отличие от фильтра нижних частот, который подавляет все частоты выше частоты среза, режекторный фильтр предназначен для подавления сигнала в узкой полосе вокруг некоторой частоты $\omega_0$, пропуская при этом как более низкие, так и более высокие частоты.

Согласно заданию, мы рассматриваем динамический фильтр второго порядка со следующей передаточной функцией:

\begin{equation*}
    W_2(p) = \frac{p^2 + a_1p + a_2}{p^2 + b_1p + b_2}
\end{equation*}

где $a_1, a_2, b_1, b_2$ являются вещественными параметрами, которые нам предстоит определить. Для их нахождения воспользуемся условиями, приведенными в задании.


\subsection{Определение параметров фильтра}

Получим общее выражение для амплитудно частотной характеристики фильтра.

\paragraph{Общий вид АЧХ}
\begin{equation*}
    W_2(i\omega) = \frac{(i\omega)^2 + a_1(i\omega) + a_2}{(i\omega)^2 + b_1(i\omega) + b_2} = \frac{(a_2 - \omega^2) + i(a_1\omega)}{(b_2 - \omega^2) + i(b_1\omega)}
\end{equation*}
\begin{equation*}
    |W_2(i\omega)| = \frac{|(a_2 - \omega^2) + i(a_1\omega)|}{|(b_2 - \omega^2) + i(b_1\omega)|}
\end{equation*}
\begin{equation*}
    |W_2(i\omega)| = \frac{\sqrt{(a_2 - \omega^2)^2 + (a_1\omega)^2}}{\sqrt{(b_2 - \omega^2)^2 + (b_1\omega)^2}}
\end{equation*}

\paragraph{1. Условие равенства АЧХ нулю}
Согласно заданию, для некоторой частоты $\omega_0$ АЧХ фильтра должна быть равна нулю. Дробь АЧХ фильтра равна нулю только тогда, когда её числитель равен нулю:
\begin{equation*}
    \sqrt{(a_2 - \omega_0^2)^2 + (a_1\omega_0)^2} = 0
\end{equation*}
Корень из суммы двух неотрицательных слагаемых равен нулю только в том случае, если каждое слагаемое равно нулю. Это дает нам систему из двух уравнений:
\begin{itemize}
    \item $(a_2 - \omega_0^2)^2 = 0 \implies a_2 = \omega_0^2$
    \item $(a_1\omega_0)^2 = 0 \implies a_1 = 0$ (поскольку $\omega_0$ это конкретная частота, она не равна нулю).
\end{itemize}
Таким образом, мы определили два параметра числителя: $a_1 = 0$ и $a_2 = \omega_0^2$.

\paragraph{2. Условие для крайних частот}
Фильтр должен пропускать сигналы на очень низких и очень высоких частотах, то есть его АЧХ должна стремиться к единице при $\omega \to 0$ и $\omega \to \infty$. Подставим уже найденные $a_1=0$ и $a_2=\omega_0^2$ в общую формулу АЧХ:
\begin{equation*}
    |W_2(i\omega)| = \frac{\sqrt{(\omega_0^2 - \omega^2)^2}}{\sqrt{(b_2 - \omega^2)^2 + (b_1\omega)^2}} = \frac{|\omega_0^2 - \omega^2|}{\sqrt{(b_2 - \omega^2)^2 + (b_1\omega)^2}}
\end{equation*}
Теперь проверим пределы. При $\omega \to 0$:
\begin{equation*}
    \lim_{\omega \to 0} |W_2(i\omega)| = \frac{|\omega_0^2 - 0^2|}{\sqrt{(b_2 - 0^2)^2 + (b_1 \cdot 0)^2}} = \frac{\omega_0^2}{\sqrt{b_2^2}} = \frac{\omega_0^2}{|b_2|}
\end{equation*}
По условию этот предел должен быть равен 1, следовательно, $|b_2| = \omega_0^2$. Выберем положительное значение $b_2 = \omega_0^2$.
Далее проверим предел при $\omega \to \infty$.
\begin{equation*}
    \lim_{\omega \to \infty} \frac{|\omega_0^2 - \omega^2|}{\sqrt{(\omega^2(\frac{b_2}{\omega^2} - 1))^2 + \omega^2 b_1^2}} = \lim_{\omega \to \infty} \frac{\omega^2|\frac{\omega_0^2}{\omega^2} - 1|}{\omega^2\sqrt{(\frac{b_2}{\omega^2} - 1)^2 + \frac{b_1^2}{\omega^2}}} = \frac{|0-1|}{\sqrt{(0-1)^2+0}} = 1
\end{equation*}
Этот предел равен единице и выполняется автоматически при найденных параметрах.

\paragraph{3. Условие устойчивости}
Фильтр должен быть устойчивым. Это означает, что корни полинома знаменателя $p^2 + b_1p + b_2$ должны иметь строго отрицательную вещественную часть. Это эквивалентно требованию положительности всех его коэффициентов. Мы уже нашли, что $b_2 = \omega_0^2$, и так как $\omega_0$ это реальная частота, то $\omega_0^2 > 0$. Следовательно, для устойчивости остается единственное условие: $b_1 > 0$.

\paragraph{Итоговая передаточная функция}

Собрав все найденные соотношения ($a_1=0, a_2=\omega_0^2, b_2=\omega_0^2, b_1>0$), мы получаем итоговый вид передаточной функции режекторного фильтра:

\begin{equation*}
    W_2(p) = \frac{p^2 + \omega_0^2}{p^2 + b_1p + \omega_0^2}
\end{equation*}
Здесь $\omega_0$ это частота, которую фильтр будет подавлять, а параметр $b_1 > 0$ определяет ширину и качество этой режекции. Для дальнейших исследований выберем начальное значение $b_1 = 1$.

